\chapter{Аналитический раздел}

\section{Алгоритм Хаффмана}

Алгоритм Хаффмана (англ. Huffman's algorithm)~---~алгоритм оптимального префиксного кодирования алфавита.
Был разработан в 1952 году аспирантом Массачусетского технологического института Дэвидом Хаффманом при написании им курсовой работы.
Используется во многих программах сжатия данных, например, PKZIP 2, LZH и др.

Построение кода Хаффмана сводится к построению соответствующего бинарного дерева по следующему алгоритму:
\begin{enumerate}
	\item составляется список кодируемых символов, при этом рассматривается один символ как дерево, состоящее из одного элемента c весом, равным частоте появления символа в строке;
	\item из списка выбираются два узла с наименьшим весом;
	\item формируется новый узел с весом, равным сумме весов выбранных узлов;
	\item к нему присоединяются два выбранных узла в качестве детей;
	\item к списку добавляется только что сформированный узел вместо двух объединенных узлов;
	\item если в списке больше одного узла, то пункты со второго по пятый повторяются.
\end{enumerate}